\section*{Capítulo 1}

\solucion{1.1} La moneda funcional de PETROPERÚ es el dólar: en dólares se fijan sus compras y, vía la paridad de importación, sus precios de venta. Sus estados financieros miden todo en dólares. Una deuda en soles es, por tanto, una deuda en moneda extranjera: su valor en dólares es $\text{Deuda en S/}/S$, donde $S$ es el tipo de cambio en soles por dólar. Si el sol se aprecia ($S$ baja), la deuda vale más dólares y la empresa registra una pérdida, aunque no haya pedido un sol más.

\solucion{1.2} $33.3/14.4=2.31$. Una caída de 20\,\% del WTI sobre el inventario de crudo borraría más de dos veces toda la utilidad bruta del año. El riesgo de precio es enorme en relación con la capacidad del negocio de absorberlo.

\solucion{1.3} Antes: $3{,}400/3.40=\text{US\$}~1{,}000.00$. Después: $3{,}400/3.70=\text{US\$}~918.92$. La deuda vale US\$ 81.08 menos: el deudor gana, porque el sol se depreció (cada dólar compra más soles). Es la situación simétrica a la pérdida descrita en el capítulo.

\solucion{1.4} (a) Crédito. (b) Mercado (precio de materias primas). (c) Liquidez. (d) Mercado (tasa de interés).

\solucion{1.5} Sin cobertura, en el escenario malo faltan $100-60=40$, con un costo de $0.30\times40=12$; en el bueno no hay costo. Costo esperado: $0.5\times12=6$. Con cobertura la caja es $0.5\times140+0.5\times60=100$ en ambos escenarios (suponiendo que el contrato fija el valor esperado) y no hay que endeudarse. Ahorro esperado: 6. Cuanto más volátil la caja y más caro el financiamiento externo, más vale la cobertura.

\solucion{1.6} Porque una empresa sin caja no puede pagar primas de opciones ni depositar márgenes diarios en una bolsa. Si eligiera instrumentos con esas exigencias, la propia cobertura podría provocar la crisis de liquidez que pretende evitar. Por eso el informe prefiere instrumentos OTC de costo inicial cero y con umbrales de colateral altos.
